(1) 自旋一的 $S_z$ 本征值为 $m\hbar$，$m=1,0,-1$。哈密顿量 $H=-\gamma BS_z$ 因而对应三条能级 $E_m=-\gamma B\hbar m$，即 $-\gamma B\hbar,0,+\gamma B\hbar$，相邻能级间距相同，为 $\gamma B\hbar$。
(2) $|m_x=+1\rangle$ 在 $S_z$ 基可以写成 $(|+1_z\rangle+\sqrt2|0_z\rangle+|-1_z\rangle)/2$，其分量幅度为 $1/2,1/\sqrt2,1/2$，故测得 $m_z=+1,0,-1$ 的概率分别为 $1/4,1/2,1/4$，总和严格为一。
(3) 在Hamilton量 $-\gamma BS_z$ 下，Heisenberg方程 $dS_x/dt=(i/\hbar)[H,S_x]=\gamma B S_y$、$dS_y/dt=-\gamma B S_x$。以初态 $\langle S_x(0)\rangle=\hbar$、$\langle S_y(0)\rangle=0$，解得 $\langle S_x(t)\rangle=\hbar\cos(\gamma Bt)$、$\langle S_y(t)\rangle=-\hbar\sin(\gamma Bt)$，呈Larmor进动。
(4) 非均匀磁场给不同 $m_z$ 的磁矩分量不同的力，理想Stern–Gerlach实验因此可分成三个空间束，而自旋半体系只有两个。当 $B\to0$，三Zeeman能级重新简并，时间进动频率也趋零；测量概率仍由态及测量方向确定。
